\documentclass[11pt]{article}

\usepackage[T1]{fontenc}
\usepackage[utf8]{inputenc}
\usepackage[spanish,es-nodecimaldot]{babel}
\usepackage{lmodern}
\usepackage{amsmath,amssymb,mathtools}
\usepackage{booktabs,longtable,array}
\usepackage{geometry}
\usepackage{enumitem}
\usepackage{xcolor}
\usepackage{hyperref}
\usepackage{microtype}

\geometry{margin=2.5cm}
\hypersetup{colorlinks=true,linkcolor=blue!55!black,urlcolor=blue!55!black}
\setlist[itemize]{leftmargin=1.5em,itemsep=2pt,topsep=3pt}
\setlist[enumerate]{leftmargin=1.7em,itemsep=3pt,topsep=3pt}
\renewcommand{\arraystretch}{1.18}

\newcommand{\Treat}{\mathrm{Treat}}
\newcommand{\Post}{\mathrm{Post}}
\newcommand{\ATT}{\mathrm{ATT}}
\newcommand{\SDID}{\mathrm{SDID}}
\newcommand{\FE}{\mathrm{FE}}
\newcommand{\PS}{\mathrm{PS}}
\newcommand{\ind}{\mathbb{1}}
\newcommand{\E}{\mathbb{E}}
\newcommand{\Var}{\operatorname{Var}}
\newcommand{\asinh}{\operatorname{asinh}}
\newcommand{\file}[1]{\texttt{#1}}

\title{Ecuaciones de los modelos y tablas del proyecto Free College}
\author{Apéndice metodológico autónomo}
\date{13 de octubre de 2026}

\begin{document}
\maketitle
\tableofcontents
\bigskip

\section{Objetivo y alcance}

Este documento vincula las tablas principales descritas en
\file{INDICE\_ARCHIVOS\_FREE\_COLLEGE.txt} con las ecuaciones que generan sus
estimandos. La unidad de observación es la universidad $i$ en el año $t$. El
período de análisis financiero es 2011--2022 y la reforma comienza en 2016. Los
modelos estáticos construyen controles como cambios respecto de 2015; el event
study dinámico omite 2013 para estimar explícitamente 2014 y 2015 como períodos
de anticipación.

La cohorte tratada está formada por universidades que ingresan a Gratuidad en
2016. Los controles son exclusivamente universidades nunca tratadas; las
instituciones tratadas después de 2016 no se usan como controles. Universidad
de O'Higgins y Universidad de Aysén se excluyen de la muestra financiera por no
disponer de una línea de base comparable. Las tablas principales reportan además
dos muestras lado a lado: una que conserva todos los controles disponibles y
otra que excluye UNIACC, Universidad del Pacífico, UCINF, ARCIS y Universidad de
Los Leones por cierre o pérdida de acreditación, discontinuidades de unidad o
márgenes extremos.

Definimos
\begin{align}
  \Treat_i &= \ind\{G_i=2016\}, \\
  \Post_t  &= \ind\{t\geq 2016\}, \\
  D_{it}   &= \Treat_i\Post_t,
\end{align}
donde $G_i$ es el año de entrada de la universidad $i$ a Gratuidad. El
parámetro de interés se interpreta como efecto sobre las universidades tratadas
en 2016 respecto de las nunca tratadas, bajo los supuestos de identificación de
cada método.

\section{Construcción de outcomes y transformaciones}

\subsection{Montos reales por estudiante}

Sea $P_t$ el índice de precios normalizado a $P_{2015}=1$ e $\pi_t$ la inflación
anual expresada en decimales. El índice se construye como
\begin{align}
 P_t &= P_{t-1}(1+\pi_t), && t>2015, \\
 P_t &= \frac{P_{t+1}}{1+\pi_{t+1}}, && t<2015.
\end{align}

Para una variable monetaria nominal por estudiante $M^{N}_{it}$, el nivel real
es
\begin{equation}
 M^{R}_{it}=\frac{M^{N}_{it}}{P_t}.
\end{equation}
Esta regla genera \emph{Revenue per Student Real}, \emph{OPEX per Student Real}
y \emph{CapEx per Student Real} a partir de sus variables nominales por
matrícula total.

La variación logarítmica real anual satisface
\begin{equation}
 \Delta\log M^{R}_{it}
 =\Delta\log M^{N}_{it}-\log(1+\pi_t),
\end{equation}
siempre que los valores nominales sean estrictamente positivos y que $t-1$ y
$t$ sean años consecutivos.

\subsection{Transformaciones usadas en las tablas}

Para cada outcome $j$, las tablas transformadas aplican una función $g_j(\cdot)$:
\begin{equation}
 Y^{(g)}_{it,j}=g_j(Y_{it,j}),
\end{equation}
con
\begin{align}
 g_j(y)&=\log(y), && y>0, \\
 g_j(y)&=\asinh(y)=\log\!\left(y+\sqrt{y^2+1}\right), \\
 g_j(y)&=\asinh(y/100), && \text{para ratios expresados en porcentaje}, \\
 g_j(y)&=y, && \text{para variaciones logarítmicas reales ya construidas}.
\end{align}
Los valores $y\leq 0$ son missing en las transformaciones logarítmicas; no se
imputan ni se corrigen. Para coeficientes en log, el cambio porcentual exacto es
\begin{equation}
 100\left[\exp(\widehat\tau)-1\right].
\end{equation}

\section{Tabla de estadísticas descriptivas}

\subsection*{Archivos asociados}
\begin{itemize}
 \item \file{Difference-in-Difference Analysis/results/Table\_Descriptive\_Statistics.xlsx}
 \item \file{Summary statistics before and after the reform/Table\_Excluded\_Universities\_Descriptive.xlsx}
\end{itemize}

Esta tabla no estima una ecuación causal. Para una variable $Y$ y un subconjunto
$\mathcal S$ (pre, post, CRUCH o privada), reporta
\begin{align}
 \bar Y_{\mathcal S}
 &=\frac{1}{n_{\mathcal S}}\sum_{(i,t)\in\mathcal S}Y_{it}, \\
 \widetilde Y_{\mathcal S}
 &=\operatorname{mediana}\{Y_{it}:(i,t)\in\mathcal S\}, \\
 s_{\mathcal S}
 &=\sqrt{\frac{1}{n_{\mathcal S}-1}
   \sum_{(i,t)\in\mathcal S}(Y_{it}-\bar Y_{\mathcal S})^2}, \\
 N_{\mathcal S}&=n_{\mathcal S}.
\end{align}
La hoja \file{Pre\_Post} utiliza $t<2016$ y $t\geq 2016$. La hoja
\file{CRUCH\_Private} usa la clasificación institucional fija del proyecto.
La tabla de universidades excluidas aplica los mismos momentos a cada
institución y variable, y agrega el primer y último año observado; es
descriptiva y no causal.

\section{Difference-in-Differences convencional}

\subsection*{Archivos asociados}
\begin{itemize}
 \item \file{Table\_DiD\_Main\_Levels.xlsx}
 \item \file{Table\_DiD\_Main\_Transformed.xlsx}
\end{itemize}

\subsection{Columna (1): Basic DiD}

La especificación canónica pre/post es
\begin{equation}
 Y_{it}=\alpha+\beta_1\Treat_i+\beta_2\Post_t
       +\tau D_{it}+\varepsilon_{it}.
 \label{eq:basic-did}
\end{equation}
El coeficiente reportado es $\widehat\tau$, asociado a
$D_{it}=\Treat_i\Post_t$.

\subsection{Columna (2): efectos fijos de universidad y año}

La especificación TWFE es
\begin{equation}
 Y_{it}=\alpha_i+\lambda_t+\tau D_{it}+\varepsilon_{it},
 \label{eq:twfe}
\end{equation}
donde $\alpha_i$ absorbe heterogeneidad permanente de cada universidad y
$\lambda_t$ absorbe shocks comunes de cada año.

\subsection{Columna (3): efectos fijos y controles}

La especificación completa es
\begin{equation}
 Y_{it}=\alpha_i+\lambda_t+\tau D_{it}
       +\boldsymbol\gamma'\mathbf X_{it}+\varepsilon_{it},
 \label{eq:twfe-controls}
\end{equation}
con
\begin{equation}
 \mathbf X_{it}=
 \begin{bmatrix}
  \Delta_{2015}\log(\text{matrícula}_{it}) \\
  \Delta_{2015}(\text{elegibilidad proxy}_{it}/10) \\
  \Delta_{2015}(\text{años de acreditación}_{it}) \\
  \Delta_{2015}(\text{cobertura elegibilidad}_{it}/10)
 \end{bmatrix},
 \qquad
 \Delta_{2015}X_{it}=X_{it}-X_{i,2015}.
\end{equation}

En las tres columnas, la inferencia agrupa los errores estándar por universidad:
\begin{equation}
 \widehat{\Var}_{\mathrm{cl},i}(\widehat{\boldsymbol\beta})
 =(X'X)^{-1}
 \left(\sum_i X_i'\widehat{\mathbf u}_i
              \widehat{\mathbf u}_i'X_i\right)
 (X'X)^{-1},
\end{equation}
con la corrección finita implementada por el software.

La tabla transformada reemplaza $Y_{it}$ por $g_j(Y_{it})$ en
\eqref{eq:basic-did}--\eqref{eq:twfe-controls}; no cambia el diseño.

\subsection{Event study dinámico}

Los archivos \file{event\_dynamic\_<OUTCOME>.csv} y
\file{event\_log\_dynamic\_<OUTCOME>.csv} corresponden a
\begin{equation}
 Y_{it}=\alpha_i+\lambda_t+
 \sum_{k\neq -3}\tau_k
 \left[\Treat_i\ind\{t-2016=k\}\right]
 +\boldsymbol\gamma'\mathbf X_{it}+\varepsilon_{it},
 \label{eq:event-study}
\end{equation}
donde $k=-3$ (año 2013) es la categoría omitida. El coeficiente $k=-4$ (2012)
es el único coeficiente estrictamente anterior a la referencia disponible en la
ventana estimada; $k=-2$ y $k=-1$ corresponden a 2014 y 2015 y se interpretan
como anticipación. Los coeficientes $k\geq 0$ representan la dinámica posterior
a la reforma bajo el supuesto de tendencias paralelas.

\section{Pre-tendencias, anticipación y sensibilidad de Roth}

\subsection*{Archivos asociados}
\begin{itemize}
 \item \file{Pre-Trends Analysis/Results/Table\_PreTrends\_Main\_Levels.xlsx}
 \item \file{Pre-Trends Analysis/Results/Table\_PreTrends\_Main\_Transformed.xlsx}
 \item \file{Pre-Trends Analysis/Results/Table\_PreTrends\_Power.csv}
 \item \file{Pre-Trends Analysis/Results/Table\_HonestDiD\_Relative\_Magnitude.csv}
 \item \file{Pre-Trends Analysis/Results/Table\_HonestDiD\_Smoothness.csv}
\end{itemize}

Este módulo es independiente de las tablas DiD, SDID e IV. Las dos tablas
principales estiman, para $h\in\{2012,2014,2015\}$, la misma regresión de
cambios usada para construir el event study DiD:
\begin{equation}
 Y_{ih}-Y_{i,2013}
 =a_h+\theta_h\Treat_i
 +\boldsymbol\gamma_h'
   (\mathbf X_{ih}-\mathbf X_{i,2013})+u_{ih},
 \label{eq:pretrends-change-regression}
\end{equation}
donde $\mathbf X_{it}$ contiene log matrícula de pregrado, proxy de elegibilidad
dividida por 10, años de acreditación y cobertura del proxy dividida por 10.
Así, $\theta_{2012}$ es la pre-tendencia reportada y $\theta_{2014}$ y
$\theta_{2015}$ son los coeficientes de anticipación. Cada regresión se estima
por separado para la muestra con y sin los cinco controles de riesgo. Los
errores estándar y los IC al 95\% provienen de bootstrap por universidad; el
$R^2$ ajustado corresponde a la regresión OLS de cambios de
\eqref{eq:pretrends-change-regression}. Para la tabla transformada se reemplaza
$Y$ por $g_j(Y)$, usando la misma transformación que en DiD y SDID.

Los diagnósticos \file{pretrends} y \file{HonestDiD} reciben del event study el
vector de coeficientes y su matriz de varianza--covarianza conjunta,
\begin{equation}
 \widehat{\boldsymbol\beta}
 =\left(\widehat\tau_k:k\in\mathcal K,\ k\neq-3\right)',
 \qquad
 \widehat{\boldsymbol\Sigma}
 =\widehat{\Var}_{\mathrm{bootstrap},i}
   (\widehat{\boldsymbol\beta}),
 \label{eq:roth-inputs}
\end{equation}
donde el bootstrap remuestrea universidades completas para preservar la
dependencia serial dentro de institución.

\subsection{Potencia del test de pre-tendencias}

El paquete \file{pretrends} evalúa una familia de violaciones lineales respecto
de la referencia 2013,
\begin{equation}
 \delta_k(s)=s\{k-(-3)\}=s(k+3),
 \label{eq:linear-pretrend}
\end{equation}
y calcula la pendiente mínima cuya prueba de pre-tendencias tendría potencia
objetivo de 80\%:
\begin{equation}
 s_{0.80}
 =\inf\left\{s\geq0:
 \Pr_{\boldsymbol\delta(s)}\!\left(\phi_{\mathrm{pre}}=1\right)
 \geq0.80\right\}.
 \label{eq:pretrend-power}
\end{equation}
Aquí $\phi_{\mathrm{pre}}=1$ indica que el test detecta al menos una desviación
pretratamiento significativa. Un $s_{0.80}$ grande implica baja potencia contra
tendencias lineales económicamente moderadas. Como con base 2013 solo existe un
coeficiente anterior ($k=-4$), el resultado es un diagnóstico explícito de la
limitación de potencia y no una validación de tendencias paralelas. Esta es la
implementación de \href{https://github.com/jonathandroth/pretrends}{Roth
(2022), \file{pretrends}}.

\subsection{Intervalos robustos HonestDiD}

Sea $\Delta_k$ la desviación de la tendencia contrafactual entre tratadas y
controles. El coeficiente dinámico observado puede escribirse como
\begin{equation}
 \beta_k=\tau_k+\Delta_k,
 \label{eq:event-bias-decomposition}
\end{equation}
donde $\tau_k$ es el efecto causal y $\Delta_k$ es el sesgo por violación de
tendencias paralelas. El análisis de magnitud relativa restringe los cambios
post-referencia mediante
\begin{equation}
 |\Delta_{k+1}-\Delta_k|
 \leq \bar M
 \max_{\ell<-3}|\Delta_{\ell+1}-\Delta_\ell|,
 \qquad k\geq-3.
 \label{eq:honest-relative}
\end{equation}
Por ejemplo, $\bar M=1$ permite que la peor desviación posterior sea tan grande
como la peor variación observada antes de la referencia; $\bar M=2$ permite el
doble. La restricción de suavidad acota cambios en la pendiente:
\begin{equation}
 |(\Delta_{k+1}-\Delta_k)-(\Delta_k-\Delta_{k-1})|\leq M.
 \label{eq:honest-smoothness}
\end{equation}
Con $M=0$ la diferencia contrafactual de tendencias es lineal; valores mayores
permiten más curvatura.

Para el outcome priorizado, crecimiento real anual del ingreso operacional por
estudiante, se reportan dos funcionales:
\begin{align}
 \theta_{2016}&=\tau_0, \\
 \theta_{2016:2021}&=\frac{1}{6}\sum_{k=0}^{5}\tau_k.
 \label{eq:honest-estimands}
\end{align}
Si $\boldsymbol\ell$ selecciona uno de estos funcionales y
$\mathcal D(c)$ representa el conjunto de desviaciones admisibles bajo
\eqref{eq:honest-relative} o \eqref{eq:honest-smoothness}, HonestDiD construye
un conjunto de confianza uniforme sobre esas desviaciones:
\begin{equation}
 \mathrm{CS}_{1-\alpha}(c)
 =\bigcup_{\boldsymbol\Delta\in\mathcal D(c)}
 \mathrm{CI}_{1-\alpha}
 \left[\boldsymbol\ell'
 (\widehat{\boldsymbol\beta}-\boldsymbol\Delta);
 \widehat{\boldsymbol\Sigma}\right].
 \label{eq:honest-cs}
\end{equation}
El código usa intervalos condicionales de longitud fija para magnitud relativa y
FLCI para suavidad, siguiendo
\href{https://github.com/asheshrambachan/HonestDiD}{Rambachan y Roth (2023),
\file{HonestDiD}}. Un punto sin solución numérica se reporta como ``--'' y no se
interpreta como evidencia a favor o en contra del efecto.

\section{Synthetic Difference-in-Differences}

\subsection*{Archivos asociados}
\begin{itemize}
 \item \file{Table\_SDID\_Main\_Levels.xlsx}
 \item \file{Table\_SDID\_Main\_Transformed.xlsx}
 \item \file{results/SDID/SDID\_Diagnostics.xlsx}
\end{itemize}

Para cada outcome se construye una matriz balanceada $Y\in\mathbb R^{N\times T}$.
Las primeras $N_0$ filas son controles y las últimas $N_1$ filas son tratadas;
las primeras $T_0$ columnas son pretratamiento y las últimas $T_1$ son
postratamiento.

\subsection{Pesos de universidades y de períodos}

Los pesos donor $\boldsymbol\omega=(\omega_1,\ldots,\omega_{N_0})'$ satisfacen
\begin{equation}
 \omega_i\geq 0,\qquad \sum_{i=1}^{N_0}\omega_i=1,
\end{equation}
y son elegidos para aproximar la trayectoria pretratamiento promedio de las
tratadas mediante una combinación ponderada de controles, con regularización.
En notación compacta,
\begin{equation}
 (\widehat\omega_0,\widehat{\boldsymbol\omega})
 =\arg\min_{\omega_0,\boldsymbol\omega\in\Omega}
 \left\|
 \overline{\mathbf Y}_{1,\mathrm{pre}}
 -\omega_0\mathbf 1
 -\mathbf Y_{0,\mathrm{pre}}'\boldsymbol\omega
 \right\|_2^2
 +\zeta_\omega^2\|\boldsymbol\omega\|_2^2.
 \label{eq:omega}
\end{equation}

Los pesos temporales $\boldsymbol\lambda=(\lambda_1,\ldots,\lambda_{T_0})'$,
con $\lambda_t\geq0$ y $\sum_{t<T_0}\lambda_t=1$, se eligen para balancear
las diferencias entre controles en el pretratamiento y el postratamiento:
\begin{equation}
 (\widehat\lambda_0,\widehat{\boldsymbol\lambda})
 =\arg\min_{\lambda_0,\boldsymbol\lambda\in\Lambda}
 \left\|
 \overline{\mathbf Y}_{0,\mathrm{post}}
 -\lambda_0\mathbf 1
 -\mathbf Y_{0,\mathrm{pre}}\boldsymbol\lambda
 \right\|_2^2
 +\zeta_\lambda^2\|\boldsymbol\lambda\|_2^2.
 \label{eq:lambda}
\end{equation}

\subsection{Estimador y curva de efectos}

Definimos la diferencia tratada--sintética no ajustada en cada año:
\begin{equation}
 \delta_t=
 \frac{1}{N_1}\sum_{i=N_0+1}^{N}Y_{it}
 -\sum_{i=1}^{N_0}\widehat\omega_iY_{it}.
\end{equation}
El offset pretratamiento es
\begin{equation}
 \widehat b_{\mathrm{pre}}
 =\sum_{t=1}^{T_0}\widehat\lambda_t\delta_t.
\end{equation}
Por tanto, la curva postratamiento es
\begin{equation}
 \widehat\tau_t^{\SDID}=\delta_t-\widehat b_{\mathrm{pre}},
 \qquad t>T_0,
 \label{eq:sdid-curve}
\end{equation}
y el ATT principal reportado en la tabla es
\begin{equation}
 \widehat\tau^{\SDID}
 =\frac{1}{T_1}\sum_{t=T_0+1}^{T}
   \widehat\tau_t^{\SDID}.
 \label{eq:sdid-att}
\end{equation}

Equivalentemente, $\widehat\tau^{\SDID}$ es el coeficiente de $D_{it}$ en la
regresión ponderada
\begin{equation}
 (\widehat\mu,\widehat{\boldsymbol\alpha},
  \widehat{\boldsymbol\lambda}^{\FE},\widehat\tau)
 =\arg\min_{\mu,\boldsymbol\alpha,\boldsymbol\lambda^{\FE},\tau}
 \sum_{i,t}\widehat\omega_i^{*}\widehat\lambda_t^{*}
 \left(Y_{it}-\mu-\alpha_i-\lambda_t^{\FE}-\tau D_{it}\right)^2,
\end{equation}
donde los pesos tratados y postratamiento son uniformes y los pesos de controles
y del pretratamiento provienen de \eqref{eq:omega} y \eqref{eq:lambda}.

En los archivos dinámicos SDID, los valores pre-2016 son brechas de ajuste
\begin{equation}
 \widehat g_t^{\mathrm{pre}}=\delta_t-\widehat b_{\mathrm{pre}},
 \qquad t<2016,
\end{equation}
no efectos causales anticipados. Para $t\geq2016$ se usa
\eqref{eq:sdid-curve}.

La inferencia principal remuestrea universidades con reemplazo $B=500$ veces:
\begin{align}
 \widehat{\mathrm{se}}_{\mathrm{boot}}(\widehat\tau)
 &=\sqrt{\frac{B-1}{B}}
   \operatorname{sd}\!\left(
    \widehat\tau^{*(1)},\ldots,\widehat\tau^{*(B)}
   \right), \\
 \mathrm{IC}_{95\%}
 &=\widehat\tau\pm 1.96\,
   \widehat{\mathrm{se}}_{\mathrm{boot}}(\widehat\tau).
\end{align}
Si el bootstrap no entrega un error válido, se usa jackknife. El $R^2$ ajustado
no está definido para este estimador y por eso se reporta como ``--''.

\section{Propensity Score Matching y entropy balancing}

\subsection*{Archivos asociados}
\begin{itemize}
 \item \file{Propensity Score Matching/Results/Table\_PSM\_Design.csv}
 \item \file{Propensity Score Matching/Results/PSM\_weights.csv}
\end{itemize}

Estos archivos documentan el diseño de comparabilidad pretratamiento; no
contienen por sí solos una regresión de outcomes. El propensity score se estima
mediante logit:
\begin{equation}
 e(\mathbf X_i)
 =\Pr(\Treat_i=1\mid\mathbf X_i)
 =\Lambda(\alpha+\mathbf X_i'\boldsymbol\beta),
 \qquad
 \Lambda(z)=\frac{e^z}{1+e^z}.
\end{equation}

En el diseño principal,
\begin{equation}
 \mathbf X_i=
 \begin{bmatrix}
  \log(\text{matrícula}_{i,2015}) \\
  \log(\text{arancel}_{i,2015}) \\
  \text{Asset Turnover}_{i,2015}
 \end{bmatrix}.
\end{equation}
La distancia de Mahalanobis entre la tratada $i$ y el control $j$ es
\begin{equation}
 d_M(i,j)=
 \sqrt{(\mathbf X_i-\mathbf X_j)'
 S_X^{-1}(\mathbf X_i-\mathbf X_j)},
\end{equation}
sujeta al caliper de propensity score
\begin{equation}
 \left|\operatorname{logit}\widehat e_i-
       \operatorname{logit}\widehat e_j\right|
 \leq 0.20\,
 \operatorname{sd}\!\left(\operatorname{logit}\widehat e\right).
\end{equation}
Se seleccionan dos controles por tratada, con reemplazo. El peso de un control
es proporcional a su frecuencia de uso como match; las tratadas reciben peso
uno.

El balance de la covariable $X_k$ se resume mediante
\begin{equation}
 \mathrm{SMD}_k=
 \frac{\bar X_{1k}-\bar X_{0k,w}}
 {\sqrt{(s_{1k}^2+s_{0k,w}^2)/2}},
\end{equation}
y el tamaño muestral efectivo de los controles es
\begin{equation}
 \mathrm{ESS}_0=\frac{\left(\sum_{i:\Treat_i=0}w_i\right)^2}
 {\sum_{i:\Treat_i=0}w_i^2}.
\end{equation}

En la columna de entropy balancing, los pesos resuelven
\begin{equation}
 \min_{w_i>0}\sum_{i:\Treat_i=0}w_i\log\left(\frac{w_i}{q_i}\right)
\end{equation}
sujeto a
\begin{equation}
 \sum_{i:\Treat_i=0}w_i=1,
 \qquad
 \sum_{i:\Treat_i=0}w_i\mathbf X_i
 =\frac{1}{N_1}\sum_{i:\Treat_i=1}\mathbf X_i.
\end{equation}

Si posteriormente se estima un ATT con estos pesos, el contraste básico es
\begin{equation}
 \widehat\tau_{\ATT,w}
 =\frac{1}{N_1}\sum_{i:\Treat_i=1}Y_i
 -\sum_{i:\Treat_i=0}\widetilde w_iY_i,
 \qquad \sum_{i:\Treat_i=0}\widetilde w_i=1.
\end{equation}
No obstante, ese contraste no aparece en \file{Table\_PSM\_Design.csv}: dicha
tabla evalúa diseño, retención, soporte, balance y ESS.

\section{Panel de variables instrumentales}

\subsection*{Archivos asociados}
\begin{itemize}
 \item Código fuente: \file{Instrumental Variables Analysis/code/IV\_fixest\_FreeCollege.R}.
 \item Resultados principales: \file{results\_ols\_iv\_fc5.csv},
       \file{results\_baseline\_iv\_fc5.csv} y
       \file{results\_conditional\_iv\_fc5.csv}.
 \item Primer estadio y diagnósticos: \file{results\_first\_stage\_fc5.csv}
       y \file{diagnostics\_iv\_fc5.csv}.
 \item Robusteces: \file{results\_iv\_robustness\_fc5.csv},
       \file{results\_reduced\_form\_fc5.csv},
       \file{results\_pre2016\_placebos\_fc5.csv} y los archivos relativos
       al instrumento de arancel de referencia.
\end{itemize}

Esta sección reproduce el código vigente y no la parametrización histórica de
cuatro instrumentos que aparecía en una versión anterior de la tabla. El
tratamiento endógeno es la participación observada y absorbente en Gratuidad,
$D_{it}=\Treat_i\Post_t$, para la cohorte de adopción de 2016 frente a
universidades nunca tratadas.

Sean
\begin{align}
 A_{i,2015}&=\text{años de acreditación observados en 2015}, \\
 E_{i,2015}&=\text{elegibilidad proxy de 2015, dividida por 10}, \\
 T_{i,2006}^{z}&=\text{cobertura estandarizada del arancel de referencia de 2006}.
\end{align}
El instrumento principal del código es
\begin{align}
 Z^{A4}_{it}
 &=\ind\{A_{i,2015}\geq4\}\Post_t.
 \label{eq:iv-main-instrument}
\end{align}
Es una característica predeterminada de la institución, fijada antes de la
reforma, e interactuada con el período post-2016.

\subsection{Referencia OLS con efectos fijos}

La referencia no instrumentada es
\begin{equation}
 Y_{it}=\alpha_i+\lambda_t+\tau_{\mathrm{OLS}}D_{it}+u_{it}.
 \label{eq:iv-ols}
\end{equation}
Esta especificación no incorpora controles contemporáneos, pues matrícula,
elegibilidad y acreditación podrían ser mediadores de la política. Por tanto,
$\tau_{\mathrm{OLS}}$ se interpreta como el efecto total asociativo dentro de
la comparación con efectos fijos.

\subsection{IV-FE principal: primer estadio y ecuación estructural}

El modelo principal es exactamente identificado. Su primer estadio es
\begin{equation}
 D_{it}=a_i+b_t+\pi_A Z^{A4}_{it}+v_{it},
 \label{eq:first-stage}
\end{equation}
y la ecuación estructural estimada por 2SLS es
\begin{equation}
 Y_{it}=\alpha_i+\lambda_t+\tau_{\mathrm{LATE}}\widehat D_{it}+u_{it}.
 \label{eq:second-stage}
\end{equation}
Aquí $\widehat D_{it}$ proviene de \eqref{eq:first-stage}. Bajo relevancia,
independencia, exclusión y monotonicidad,
$\tau_{\mathrm{LATE}}$ identifica un efecto local de la participación en
Gratuidad para las instituciones cuyo tratamiento responde al umbral de
acreditación. Los errores estándar se agrupan por universidad.

\subsection{Especificaciones condicionales}

El código estima dos variantes secundarias con controles contemporáneos. Si
$X_{it}$ denota uno de los siguientes vectores,
\begin{align}
 X^{(c)}_{it}
 &=\bigl(\log \text{matrícula}_{it},E_{it}/10,C_{it}/10,A_{it}\bigr)', \\
 X^{(d)}_{it}
 &=\bigl(\log \text{matrícula}_{it},E_{it}/10,C_{it}/10,
          \ind\{A_{it}\geq4\}\bigr)',
\end{align}
donde $C_{it}$ es la cobertura contemporánea de la elegibilidad proxy, las
ecuaciones son
\begin{align}
 D_{it}
 &=a_i+b_t+X_{it}'\gamma+\pi_A Z^{A4}_{it}+v_{it},
 \label{eq:first-stage-controls}\\
 Y_{it}
 &=\alpha_i+\lambda_t+X_{it}'\beta
   +\tau_{\mathrm{cond}}\widehat D_{it}+u_{it}.
 \label{eq:second-stage-controls}
\end{align}
Como esos controles pueden responder a Gratuidad, el código separa estas
estimaciones condicionales del efecto total principal.

\subsection{Instrumentos y especificaciones de robustez}

Además de \eqref{eq:iv-main-instrument}, el código construye
\begin{align}
 Z^{E15}_{it}&=E_{i,2015}\Post_t, \\
 Z^{A\mathrm{yrs}}_{it}&=A_{i,2015}\Post_t, \\
 Z^{T06}_{it}&=T_{i,2006}^{z}\Post_t, \\
 Z^{\bar E}_{it}&=\bar E_{i,2014:2015}\Post_t, \\
 Z^{\bar A}_{it}&=\bar A_{i,t<2016}\Post_t, \\
 Z^{E\mathrm{cur}}_{it}&=(E_{it}/10)\Post_t, \\
 Z^{A4\mathrm{cur}}_{it}&=\ind\{A_{it}\geq4\}\Post_t.
\end{align}
La cobertura de arancel se define en los datos como
\begin{equation}
 \mathrm{Coverage}_{i,2006}
 =\frac{\text{arancel de referencia}_{i,2006}}
        {\text{arancel efectivo}_{i,2006}}
 =1+\frac{\text{referencia}-\text{efectivo}}
          {\text{efectivo}},
\end{equation}
y $T_{i,2006}^{z}$ es su versión estandarizada.

Cada robustez sustituye $Z^{A4}_{it}$ en
\eqref{eq:first-stage} por el conjunto $\mathcal Z_{it}^{(r)}$:
\begin{align}
 \mathcal Z^{(1)}&=\{Z^{E15}\}, &
 \mathcal Z^{(2)}&=\{Z^{A4}\}, &
 \mathcal Z^{(3)}&=\{Z^{T06}\}, \\
 \mathcal Z^{(4)}&=\{Z^{A4},Z^{T06}\}, &
 \mathcal Z^{(5)}&=\{Z^{E15},Z^{A\mathrm{yrs}}\}, \\
 \mathcal Z^{(6)}&=\{Z^{\bar E},Z^{\bar A}\}, &
 \mathcal Z^{(7)}&=\{Z^{E\mathrm{cur}},Z^{A4\mathrm{cur}}\}.
\end{align}
El último conjunto es una robustez potencialmente endógena y no un diseño IV
principal. El código también repite el IV base en una muestra común de outcomes.

Para los dos instrumentos predeterminados destacados, las formas reducidas son
\begin{align}
 Y_{it}&=\alpha_i+\lambda_t+\rho_A Z^{A4}_{it}+e_{it}, \\
 Y_{it}&=\alpha_i+\lambda_t+\rho_T Z^{T06}_{it}+e_{it}.
\end{align}

\subsection{Placebos y diagnósticos}

Para cada falso inicio $f\in\{2013,2014,2015\}$ y usando solamente años
anteriores a 2016, el código estima
\begin{equation}
 Y_{it}=\alpha_i+\lambda_t
 +\rho_{E,f}\bigl(E_{i,2015}\ind\{t\geq f\}\bigr)
 +\rho_{A,f}\bigl(\ind\{A_{i,2015}\geq4\}\ind\{t\geq f\}\bigr)
 +\eta\log(\text{matrícula}_{it})+e_{it}.
\end{equation}

La fortaleza del primer estadio se evalúa prioritariamente con el estadístico
Wald IV robusto a la agrupación. El código también registra el estadístico $F$
del primer estadio, Wu--Hausman y, cuando existe sobreidentificación, Sargan.
Los outcomes en variaciones reales logarítmicas anuales se estiman directamente;
el código no aplica una transformación \emph{signed log1p} a niveles o razones.

\section{Correspondencia resumida entre tablas y estimandos}

\begin{longtable}{>{\raggedright\arraybackslash}p{0.37\textwidth}
                  >{\raggedright\arraybackslash}p{0.25\textwidth}
                  >{\raggedright\arraybackslash}p{0.27\textwidth}}
\toprule
Archivo & Estimando o contenido & Ecuación principal \\
\midrule
\endfirsthead
\toprule
Archivo & Estimando o contenido & Ecuación principal \\
\midrule
\endhead
\file{Table\_Descriptive\_Statistics.xlsx}
& Momentos muestrales pre/post y CRUCH/privada
& Media, mediana, desviación estándar y $N$ \\
\addlinespace
\file{Table\_Excluded\_Universities\_Descriptive.xlsx}
& Momentos de las cinco universidades excluidas
& Media, mediana, desviación estándar, extremos y cobertura temporal \\
\addlinespace
\file{Table\_DiD\_Main\_Levels.xlsx}
& Efecto promedio DiD en niveles
& Ecuaciones \eqref{eq:basic-did}--\eqref{eq:twfe-controls} \\
\addlinespace
\file{Table\_DiD\_Main\_Transformed.xlsx}
& Efecto DiD sobre $g_j(Y)$
& Ecuaciones \eqref{eq:basic-did}--\eqref{eq:twfe-controls} con outcome transformado \\
\addlinespace
\file{Table\_PreTrends\_Main\_Levels.xlsx}
& Pre-tendencia 2012 y anticipación 2014--2015 en niveles
& Ecuación \eqref{eq:pretrends-change-regression} \\
\addlinespace
\file{Table\_PreTrends\_Main\_Transformed.xlsx}
& Pre-tendencia 2012 y anticipación 2014--2015 sobre $g_j(Y)$
& Ecuación \eqref{eq:pretrends-change-regression} con outcome transformado \\
\addlinespace
\file{Table\_HonestDiD\_Relative\_Magnitude.csv} y
\file{Table\_HonestDiD\_Smoothness.csv}
& IC robustos para impacto 2016 y promedio 2016--2021
& Ecuaciones \eqref{eq:event-bias-decomposition}--\eqref{eq:honest-cs} \\
\addlinespace
\file{Table\_SDID\_Main\_Levels.xlsx}
& ATT SDID ponderado en niveles
& Ecuaciones \eqref{eq:omega}--\eqref{eq:sdid-att} \\
\addlinespace
\file{Table\_SDID\_Main\_Transformed.xlsx}
& ATT SDID ponderado sobre $g_j(Y)$
& Ecuaciones \eqref{eq:omega}--\eqref{eq:sdid-att} con outcome transformado \\
\addlinespace
\file{Table\_PSM\_Design.csv}
& Calidad del diseño de matching/ponderación
& Propensity score, Mahalanobis, SMD y ESS \\
\addlinespace
\file{PSM\_weights.csv}
& Pesos ATT de universidades
& Frecuencia de matching y entropy balancing \\
\addlinespace
\file{Table\_IV\_Main.xlsx}
& OLS-FE e IV-FE, con y sin controles de riesgo
& Ecuaciones \eqref{eq:iv-ols}--\eqref{eq:second-stage} \\
\addlinespace
\file{results\_conditional\_iv\_fc5.csv}
& IV-FE con controles contemporáneos
& Ecuaciones \eqref{eq:first-stage-controls} y \eqref{eq:second-stage-controls} \\
\addlinespace
\file{results\_iv\_robustness\_fc5.csv}
& Instrumentos alternativos y combinados
& Primera y segunda etapas con $\mathcal Z^{(r)}$ \\
\bottomrule
\end{longtable}

\section{Supuestos que deben acompañar la interpretación}

\begin{enumerate}
 \item \textbf{DiD/TWFE:} tendencias paralelas condicionales, ausencia de
 anticipación y composición comparable de la muestra.
 \item \textbf{Pretrends/HonestDiD:} el análisis de potencia no prueba
 tendencias paralelas. Los intervalos robustos son válidos condicionalmente al
 conjunto de desviaciones elegido (magnitud relativa o suavidad) y a la
 aproximación normal conjunta de los coeficientes event-study.
 \item \textbf{SDID:} la combinación ponderada de controles reproduce de forma
 suficiente la trayectoria contrafactual pretratamiento; los gaps pre son
 diagnósticos de ajuste, no ATT anticipados.
 \item \textbf{PSM:} independencia condicional dada la información
 pretratamiento y soporte común. El balance observado no prueba ausencia de
 confusión no observada.
 \item \textbf{IV:} relevancia, exclusión, independencia y monotonicidad de los
 instrumentos. Un primer estadio fuerte no valida por sí solo la restricción de
 exclusión.
\end{enumerate}

\end{document}
